% Generated draft: requires mathematical and source review.
\documentclass[11pt]{article}
\usepackage[margin=25mm]{geometry}
\usepackage{fontspec}
\setmainfont{Latin Modern Roman}
\usepackage{amsmath,amssymb,mathtools,unicode-math}
\setmathfont{Latin Modern Math}
\usepackage{enumitem,needspace,xurl,hyperref}
\setlength{\parindent}{0pt}
\setlength{\parskip}{5pt}
\setlength{\emergencystretch}{3em}
\newcommand{\N}{\mathbb{N}}
\newcommand{\Z}{\mathbb{Z}}
\newcommand{\Q}{\mathbb{Q}}
\newcommand{\R}{\mathbb{R}}
\newcommand{\C}{\mathbb{C}}
\newcommand{\F}{\mathbb{F}}
\newcommand{\A}{\mathbb{A}}
\newcommand{\PP}{\mathbb{P}}
\newcommand{\E}{\mathbb{E}}
\newcommand{\eps}{\varepsilon}
\newcommand{\dd}{\,\mathrm{d}}
\begin{document}
\section*{UnsolvedMath ID: 30002775}
\textbf{Source problem code:} catalogue entry
\subsection{Definitions and mathematical statement}
Let \(P\) be a convex simplicial polytope of dimension \(d\), and let \(\Delta=\partial P\) be its boundary simplicial complex. We use the convention \(d\geq 1\); the supplied record does not specify a lower-dimensional convention. The polytope is \emph{balanced} if there is a coloring \(c:V(\Delta)\to\{1,\ldots,d\}\) such that the endpoints of every edge have different colors. Equivalently, every facet of \(\Delta\) contains exactly one vertex of each color.

For \(-1\leq j\leq d-1\), let \(f_j(\Delta)\) be the number of \(j\)-dimensional faces, with \(f_{-1}(\Delta)=1\). Define the \(h\)-numbers \(h_0(P),\ldots,h_d(P)\) by
\[
\sum_{i=0}^{d}h_i(P)t^{d-i}
=
\sum_{i=0}^{d}f_{i-1}(\Delta)(t-1)^{d-i}.
\]
The inequality assertion in the supplied record is that every balanced simplicial \(d\)-polytope satisfies
\[
\frac{h_{r-1}(P)}{\binom{d}{r-1}}
\leq
\frac{h_r(P)}{\binom{d}{r}}
\qquad
\text{for every integer }1\leq r\leq\left\lfloor\frac d2\right\rfloor.
\]
This index range is empty when \(d=1\).

The remaining proposed target concerns equality
\[
\frac{h_{r-1}(P)}{\binom{d}{r-1}}
=
\frac{h_r(P)}{\binom{d}{r}}.
\]
The source says that equality at the relevant adjacent terms should characterize the ``balanced stacked property.'' It does not define that property, specify which indices are relevant, or state any index-dependent version or additional dimension restrictions. Consequently, a precise necessary-and-sufficient equality characterization cannot be extracted from this record without further source clarification. No particular definition of balanced stackedness or equality equivalence is imposed here.

The supplied dated curation attributes a proof of the inequality assertion to Juhnke--Kubitzke and Murai and describes the equality characterization as unresolved in the extracted context. These are supplied literature assessments, not independently verified claims.
\subsection{Short English statement}
For a balanced simplicial polytope, the supplied assertion says that its \(h\)-numbers, divided by the corresponding binomial coefficients, increase weakly through the middle index. The intended remaining question is to characterize equality using a balanced stacked property, but the record omits the definition and the relevant index restrictions. Its supplied status is partially solved.
\subsection{Sources}
\begin{itemize}
\item Supplied UnsolvedMath local registry record, permanent ID \(30002775\), title: \emph{Equality Cases in the Balanced Lower Bound Conjecture}; exact source problem code: \texttt{catalogue entry}. Collection: Oberwolfach Reports Open Problems; proposed year recorded as 2015. No individual proposer is supplied.
\item \emph{Geometric and Algebraic Combinatorics} (2015), Oberwolfach Reports, pp.\ 285--368, DOI 10.4171/owr/2015/5. Supplied report reference: \url{https://doi.org/10.4171/owr/2015/5}. The background describes these reports as records of mathematical workshops and their open-problem discussions.
\item Supplied literature source labeled ``Juhnke--Kubitzke--Murai, primary paper'': \url{https://arxiv.org/abs/1503.06430}. The supplied assessment credits Juhnke--Kubitzke and Murai with proving the inequality part of the balanced generalized lower bound conjecture.
\item Supplied literature source labeled ``Murai, Oberwolfach Report'': \url{https://doi.org/10.4171/owr/2015/5}.
\item Literature checked at: \texttt{2026-08-21}. This exact date records supplied dated curation, not independent verification. That curation labels the problem partially solved and says that the equality characterization remains unresolved in the extracted context; the background explicitly calls for independent verification of time-sensitive status and cited claims.
\end{itemize}
\end{document}
